\documentclass[10pt,a4paper]{amsart}
\usepackage[margin=19mm]{geometry}
\usepackage{amsmath,amssymb,mathtools}
\usepackage[scr=rsfs,cal=euler]{mathalfa}
\usepackage{xfrac}
\usepackage[unicode,hidelinks,bookmarks=false]{hyperref}
\newcommand{\ie}{i.e.\ }
\newcommand{\cf}{cf.\ }
\newcommand{\resp}{respectively\ }
\newcommand{\Subsection}{\subsection}
\providecommand{\nobreakdash}{\nobreak\hbox{-}}
\setlength{\emergencystretch}{2em}
\multlinegap0pt
\newtheorem{theorem}{Theorem}
\title[Edge-coloring the intersection power graph (Theorem 3.2)]{Edge-Coloring Power Graphs: the intersection power graph is Class 2 exactly for cyclic groups of odd prime-power order (Theorem 3.2)}
\author{Elie Feinsilber}
\date{Source: arXiv:2508.01417v3 (CC BY 4.0)}
\begin{document}
\maketitle

\noindent\emph{Source and licence.} Edge-Coloring Power Graphs. Version arXiv:2508.01417v3 (CC BY 4.0), distributed under the Creative Commons Attribution 4.0 International licence, as recorded in the arXiv licence metadata for this exact version and shown on the abstract page https://arxiv.org/abs/2508.01417v3. Full rights notice: licenses/arxiv-2508.01417-CC-BY-4.0.txt.\par\smallskip

\noindent\textit{Editorial scope.} Counted unit: the theorem printed as Theorem 3.2 of the article (source0.tex lines 547--551, source label \texttt{thm:intersection-power}), delivered together with the article's own complete original proof of it (source0.tex lines 553--592, one \texttt{proof} environment of 40 lines that proves the stated classification with no gap). No mathematics is written, repaired or re-derived: statement and proof are the article's own text line for line, in the article's own notation ($G$, $n$, $\mathscr{G}_{\mathrm{I}}(G)$, $\Delta$, $\chi'$, $p$, $a$, $\pi$, $H$, $K$). Required context retained uncounted: the sentence in which the article defines its intersection power graph $\mathscr{G}_{\mathrm{I}}(G)$ (lines 542--545; the retained fragment begins with the words ``denoted by'' because the immediately preceding clause of the same sentence attributes the notion to Bera by a citation command, and only that attribution clause is left out -- no mathematical content of the definition lies in it). Every symbol, condition and label of the retained lines is retained unchanged. Omitted from this package and not redistributed: the article's preamble and title page (lines 1--137); its abstract (lines 138--143); its whole introduction, including the statements of Vizing's theorem, the definitions of Class 1, Class 2 and of overfullness, the Overfull conjecture and the announcements of the article's main results (lines 144--255); its second section, where the power graph itself is edge-coloured (lines 256--338); everything in its third section other than the retained definition fragment and the counted theorem with its proof, namely the enhanced power graph theorem and its proof and the surrounding discussion (lines 339--541, 546, 593--596); its fourth section on extending colourings (lines 597--879); its fifth section on the proper power graph (lines 880--1330); its acknowledgments (lines 1331--1395); and the article's biblatex bibliography directive and closing end-document line (lines 1396--1399). Nothing inside a retained excerpt is omitted or altered: every retained body is delivered exactly as the article's lines are written, so no omission receipt of any kind accompanies this package. Citations: the retained excerpts carry no citation command at all; the article's bibliography is generated by biblatex and only the directive \texttt{\textbackslash printbibliography} is present in the source, so no bibliography entry is transcribed and no reference is redistributed. Reference adaptation: none was needed and none was made. Every reference inside the delivered unit resolves inside it: the retained statement carries the label of the counted theorem and the retained proof cites nothing, so no unresolved reference or citation is produced. Printed slips kept literal: no printed word, symbol, hypothesis or number of the counted statement or of its proof was altered. Layout and class: the article's own class is the standard LaTeX \texttt{article} class with no options, loading among others amsmath, amssymb, amsthm, mathrsfs, xy, tikz and hyperref; no publisher class or style file is used or shipped, so none travels with this package and none is redistributed. The wrapper is the lane's pinned \texttt{amsart} editorial wrapper at 10pt on A4, which declares the single theorem-like environment the retained text needs (\texttt{theorem}) and needs no macro from the article's preamble, because every control sequence in the retained spans is standard or provided by the wrapper (the calligraphic \texttt{\textbackslash mathscr} comes from the wrapper's \texttt{mathalfa} loading of rsfs). The delivered numbering is the unit's own: the counted Theorem 3.2 prints as Theorem 1. The article's own measure and page style differ from the wrapper's 10pt A4; that is a page-appearance difference of the wrapper only, and no formula, symbol, hypothesis or argument changes. Assets: no retained span contains a figure environment, a graphics-inclusion command, a PDF-page inclusion, a table or a TikZ picture; the article's own figures and drawings lie in the omitted parts of the paper, no image file is copied into this package, the package declares no asset dependency and no image file is redistributed with it. Licence: version arXiv:2508.01417v3 of this article is distributed under the Creative Commons Attribution 4.0 International licence, as recorded in the arXiv licence metadata for this exact version and shown on the abstract page https://arxiv.org/abs/2508.01417v3. A full rights notice is delivered at licenses/arxiv-2508.01417-CC-BY-4.0.txt. Characters: every retained excerpt is pure ASCII (the article's single non-ASCII character lies outside the retained spans), so the delivered engine, XeLaTeX with its default Latin Modern text font, reports no missing character. No glyph is substituted and no word is rewritten.
denoted by $\mathscr{G}_{\mathrm{I}}(G)$, has vertex set $G$. The identity is
adjacent to every other vertex, and two distinct nonidentity elements $x,y$
are adjacent if and only if
$\langle x\rangle\cap\langle y\rangle\neq\{e\}$.

\begin{theorem}
\label{thm:intersection-power}
Let $G$ be a finite group of order $n\geq 2$. Then
$\mathscr{G}_{\mathrm{I}}(G)$ is Class 2 if and only if  $G$ is cyclic of odd prime-power order. 
\end{theorem}

\begin{proof}
The identity element is adjacent to every other vertex of the
intersection power graph, so
$\Delta(\mathscr{G}_{\mathrm{I}}(G))=n-1$. Also note that if $G$ has even order, then $\mathscr{G}_{\mathrm{I}}(G)$ is Class 1. 

Suppose from now on that $n$ is
odd. If $G\cong C_{p^a}$, then any two nontrivial subgroups of $G$ have
nontrivial intersection. Therefore
$\mathscr{G}_{\mathrm{I}}(G)\cong K_n$, and its chromatic index is $n$.

Suppose next that $G$ is a noncyclic $p$-group. Since $p$ is odd, a finite
$p$-group with a unique subgroup of order $p$ is cyclic. Thus $G$ contains
two distinct subgroups $P=\langle x\rangle$ and $Q=\langle y\rangle$ of
order $p$. The vertices $x$ and $y$ are nonadjacent, and no nonidentity
element $z$ can be adjacent to both. Indeed, this would imply
$P,Q\leq\langle z\rangle$, which is impossible because a cyclic $p$-group
has a unique subgroup of order $p$.

For every $z\in G\setminus\{e,x,y\}$, at least one of $xz$ and $yz$ is
therefore a missing edge. Together with the missing edge $xy$, this gives at
least $n-2\geq(n-1)/2$ missing edges. Thus the graph is not overfull and is
of Class 1.

Finally, suppose that at least two distinct primes divide $n$. By the
Feit--Thompson theorem, $G$ is solvable. Choose a
nonempty proper subset $\pi$ of the prime divisors of $n$, and write $n=ab$,
where $a$ is the full $\pi$-part of $n$ and $b$ is the full $\pi'$-part.
Then $a,b>1$ and $\gcd(a,b)=1$. Hall's theorem gives
subgroups $H,K\leq G$ with $|H|=a$ and $|K|=b$.

Since $H\cap K=\{e\}$, no element of $H\setminus\{e\}$ is adjacent to an
element of $K\setminus\{e\}$. The graph is consequently missing at least
$(a-1)(b-1)$ edges. Since $a$ and $b$ are coprime odd integers greater than
$1$, one of them is at least $5$, and
\[
 (a-1)(b-1)-\frac{n-1}{2}
   =\frac{(a-2)(b-2)-1}{2}\geq 0.
\]
The graph is therefore not overfull and is Class 1. This completes the proof. 
\end{proof}

\end{document}
